\scanpage{54}
\uline{Intégration des algèbres $\frac12$ simpliciales}

Soit $A_*$ un groupe gradué différentiel $\frac12$ simplicial
i.e. un objet $\frac12$ simplicial dans la catégorie des groupes
gradués différentiels (: op. de dérivation de degré $+1$).
\deleted{On suppose} Donc $\forall n\geq0$
\[
A_n=A_n^*=\sum_{p\in\mathbf Z}A_n^p,
\]
avec op. différentielle
$d_n:A_n^*\mathbin{-}A_n^*$,
$A_n^p\mathbin{-}A_n^{p+1}$.
On suppose

\margin{\figure{\begin{tikzpicture}\draw (0,0)--(0,-3.2);\end{tikzpicture}}}
(1) $\forall n,\ n\geq0\quad A_n^p=0$ si $p\notin[0,n]$,
et \deleted{\unsure{le}} $A_n^*$ muni
de l'augmentation \inserted{$\unsure{\varepsilon_n}:$}
$A_0\longrightarrow A_n^*$
[\uline{NB} associée à l'unique $\sigma:\Delta_n\longrightarrow\Delta_0$;
$A_0$ est considéré comme groupe \unsure{donné de la}
graduation nulle et l'op. diff. trivial]
est une \uline{résolution} de $A_0$.

Rappelons qu'on a sur le système des $A_n$ un op.
différentiel (i.e. de degré $-1$)
\[
\partial:A_n^*\mathbin{-}A_{n-1}^*
\]
\[
\partial(\omega)=\sum_{0\leq i\leq n}\partial_{n,i}^*(\omega)(-1)^i
\]
(où les $\partial_{n,i}$, $(0\leq i\leq n)$ sont les \unsure{$n+1$ applications}
strictement croissantes $\Delta_{n-1}\mathbin{-}\Delta_n$.
On sait que $\partial(\partial\omega)=0$.

\uline{Proposition 1} Il existe un et un seul système d'hom.
de groupes
\[
A_n^*\xrightarrow{k_n}A_0
\]
satisfaisant aux conditions suivantes
\margin{\figure{\begin{tikzpicture}\draw (0,0)--(0,-2.1);\end{tikzpicture}}}
\margin{\uline{NB} le système des $k_n$ \uline{ne commute pas}
en \unsure{gén.} aux op. $\frac12$ simpliciales!}

a) (normalisation) $\quad k_0=\operatorname{id}_{A_0}$

b) (condition de degré) $\quad k_n(A_n^p)=0$ si $p\ne n$

c) (Stokes) $\quad k_n(d\omega)=k_{n-1}(\partial\omega):
\ \forall\omega\in A_n^*$.

En vertu de b), la condition c) équivaut à la même
condition, avec $\omega\in A_n^{n-1}$. On montre existence et
unicité par \inserted{\unsure{degré $\leq N$ par}} récurrence sur $N$,
en notant que pour $N=0$ c'est trivial.
Si $N>0$, \deleted{$\unsure{\bar\omega\in A}$}
il reste à démontrer ex. et un. de $k_N$, i.e. de
$k_N(\bar\omega^N)$ pour $\bar\omega^N\in A_N^N$.
Or $d\bar\omega^N=0$ puisque $A_N^{N+1}=0$
\inserted{(1)}, donc $\bar\omega^N=d\omega^{N-1}$ (1).
[En fait il suffit \inserted{ici} $H^N(A_N^*)=0$ pour $N\geq1$]
\unsure{et le} \uline{\unsure{est fixé}} par c):
\[
\unsure{k_N|k_n}(d\omega^{N-1})=\unsure{k_{N-1}|k_{n-1}}(\partial\omega)
\]
\mspage{53}
